\begin{proof}
用固定的非负径向 Schwartz 函数对行球作控制。
令 $r_s=q_s/Y'$，并在与 $S$ 互素的主元素上令
$c_\sigma(s)=1_{s\in\sigma}\chi_s(b_*)/(\tau\xi(s))$，在求逆特征之前于其他地方
以零延拓。该多项式的精确环形改写为
\[
 A_m=Y'^{-3/2}\sum_s c_\sigma(s)W_1(r_s)r_s^{-1+i\nu}
                          g_{\chi_s}(s,-m).
\]
因此展开后的平方具有公共正规化 $Y'^{-3}$，联合环形剖面包含因子
$W_1(r_{s_1})r_{s_1}^{-1+i\nu}W_1(r_{s_2})r_{s_2}^{-1-i\nu}$。
记 $s_1=Cn_1$、$s_2=Cn_2$，其中 $(n_1,n_2)=1$。对行变量作 Poisson 求和得到
因子 $Q/(q_{s_1}q_{s_2})$，而两个未正规化的 Gauss 和的完备
变换为 $q_{s_1}q_{s_2}F(s_1,s_2;j)$。因此它给出精确的前置因子
$Q/Y'^3$ 以及方程~\eqref{eq:correlation-lift} 的相关性
$F(Cn_1,Cn_2;Ck)$。
这里 $F$ 再次表示方程~\eqref{old-eq:2.11} 的完整相关性，
而不是方程~\eqref{eq:local-F} 的局部系数。
第一个 Fourier 核的宗量与 $q_Cq_k/P_a$ 同阶。

在互素的剩余上，该相关性中的因子 $L_C$ 满足
$|L_C|\le q_C$。对于随后的带符号延拓，使用其指定的
有界延拓，它在两列上都以模 $\operatorname{rad}C$ 为周期，包括 $C$ 与某一列
有公共素因子的情形。当 $C$ 的某个素因子同时整除两列时，其值
定义为零；这是对该延拓的定义，而非关于此处真实
相关性的断言。在作任何逐因子的估计之前，插入完整的
M\"obius 恒等式
\[
 1_{(n_1,n_2)=1}=\sum_{d\mid n_1,\ d\mid n_2}\mu(d)
\]
并记 $n_i=dm_i$。此时不对 $m_1,m_2$ 施加互素条件。
非零项满足 $(d,Ck)=1$。在整个这一步中使用 $\mathcal R(n_1,n_2)$ 的固定
射线延拓；零剩余特征之商将没有定义。

我们为该系数指定一个公共固定模。对 $k\ne0$
利用单位约定和引理~\ref{lem:fixed-numerator-ray} 的固定 $S$-素生成元，记
$k=\zeta k_Sk_{\mathrm{good}}$，其中
$k_S=\prod_{\mathfrak p\in S}\pi_{\mathfrak p}^{e_{\mathfrak p}}$
且 $k_{\mathrm{good}}=\prod_{p\notin S}p^{e_p}$。对于 $S$ 之外的某个主
剩余理想 $m$，互反律给出
\[
 \begin{split}
 \chi_m(k)&=\kappa_{\zeta,\boldsymbol e}(m)
       \mathcal R(k_{\mathrm{good}},m)
       \prod_{p\mid k_{\mathrm{good}}}\chi_p(m)^{e_p},\\
 \kappa_{\zeta,\boldsymbol e}(m)
       &=\chi_m(\zeta)\prod_{\mathfrak p\in S}
             \chi_m(\pi_{\mathfrak p})^{e_{\mathfrak p}\bmod6}.
 \end{split}
\]
每个显式给出的好的局部幂在公共素因子上都保留其零点，
包括 $6\mid e_p$ 的情形。引理~\ref{lem:fixed-numerator-ray}
把这有限多个 $\kappa_{\zeta,\boldsymbol e}$ 放在一个支撑于 $S$ 的固定
射线群中。此外 $\mathcal R(k_{\mathrm{good}},m)$
属于一个关于 $m$ 的固定有限族。

一次性选取 $\mathfrak l_{\rm fixed}$，使其包含主类
模与 $S$ 零掩码、$A_m$ 中每个固定射线
系数的模、$\mathcal R$ 在每个变量中的一个周期，
以及前面所有
$\kappa_{\zeta,\boldsymbol e}$ 特征的导子。该
延拓系数在主类之外，或在求逆特征之前的某个固定
非单位元处取零。该模与 $C,d,k$ 无关，且
它在 $S$ 处可能有比 $\operatorname{rad}(k)$ 更高的幂次。在应用互反律之后，
关于 $m_1,m_2$ 的联合算术
系数以如下模为周期
\[
 \mathfrak r=\operatorname{lcm}
   (\mathfrak l_{\rm fixed},\operatorname{rad}C,\operatorname{rad}(k)),
 \qquad q_{\mathfrak r}\ll_{\rm fixed}q_Cq_k.
\]
若 $k$ 在 $C$ 与 $\mathfrak l_{\rm fixed}$ 之外的一个素因子处有不可被六整除的
赋值 $e_p$，则称 $k$ 为非例外的。在该
素因子处，第一剩余列具有非主零延拓
因子 $\chi_p(m_1)^{e_p}$，至多相差一个单位标量。提升、固定
相位以及其他素因子在模 $p$ 下与 $m_1$ 无关。
它的完备均值，从而由中国剩余定理得到的完备联合均值，
为零。用 $d$ 作代换
不改变这一结论，因为 $(d,k)=1$。赋值可被六整除处的
主零掩码在例外情形中仍然存在。

频率 $k=0$ 只贡献对角项：相关性
强制 $n_1=n_2=1$，且等于 $\varphi(C)$。范数与 $Y'$
同阶的可能 $C$ 有 $O(Y')$ 个，每个满足
$\varphi(C)\le q_C\ll Y'$。因此其贡献为
$O(Q/Y')$。

对 $k\ne0$，把 Fourier 宗量分成二进壳层
$q_Cq_k/P_a\asymp R$，其中 $R=1,2,4,\ldots$，且第一个壳层
包含所有至多为二的宗量。为 $q_C\asymp C_0$ 与
$q_d\asymp D_0$ 选取二进中心 $C_0,D_0\in\{1,2,4,\ldots\}$，
其中第一个二进段包含有界的单位范围。
使用名义正剩余尺度
\[
 N=Y'/(C_0D_0),\qquad
 \mathcal Q=q_{\mathfrak r}\ll_{\rm fixed}RP_a.
\]
实际剩余范数落在 $N$ 的固定倍数之内；名义值
$N$ 允许小于一。固定任意 $A>2$ 与 $B>2$。
对每个固定的 $C,d,k$，两列的那个联合环形剖面 $W_R$
使下文所用的有限多个齐次单剖面半范数被
$O_{B,J_{\mathrm{Gr}}}((1+|\nu|)^{J_{\mathrm{Gr}}} R^{-B})$ 所界，其中固定了一个
$J_{\mathrm{Gr}}$，它是在选定这些导数阶数以及 $A,B$ 之后选取的。这
可由对第一个 Fourier 核求导得到，其 Schwartz 衰减吸收由此产生的每个
$R$ 的幂次；范数幂次的导数只付出
$1+|\nu|$ 的固定幂次。这个因子 $R^{-B}$ 保留在单剖面
Fourier 估计中，并在下文的线性格和中出现一次；
不使用小的非齐次元组半范数。

设 $a_{C,d,k}$ 为该周期算术系数。对
非例外频率，其正规化完备 Fourier 变换，
在两个剩余变量上取正规化 $\mathcal Q^{-2}$，
绝对值至多为 $q_C$，并在 $(0,0)$ 处取零。由于
$\mathcal O$ 的每个理想都是主理想，周期格具有
线性尺度 $\mathcal Q^{1/2}$。因此四维格 Poisson 求和
给出，对上述 $A$，
\begin{equation}
 \left|\sum_{\mathbf m\in\mathcal O^2}
       a_{C,d,k}(\mathbf m)W_R(\mathbf m/\sqrt N)\right|
 \ll (1+|\nu|)^{J_{\mathrm{Gr}}} R^{-B}q_C N^2
             \min\{1,(\mathcal Q/N)^A\}.
 \label{eq:gram-periodic-poisson}
\end{equation}
当 $N\ge\mathcal Q$ 时，零频率消失，其他
对偶向量的线性尺度为 $(N/\mathcal Q)^{1/2}$。阶为
$2A>4$ 的 Fourier
衰减把它们的和界定在 $(\mathcal Q/N)^A$ 之内。
当 $N<\mathcal Q$ 时，使用两个环形区域中的点计数。非空的
次单位圆环有 $N$ 被一个固定正常数从下方界定，
因此它仍包含 $O(N)$ 个格点；低于该常数时它
是空的。这就在对 $a_{C,d,k}$ 不作任何本原性或张量积假设的情况下
证明了方程~\eqref{eq:gram-periodic-poisson}。

令 $K_R=RP_a/C_0$。固定的壳层比较常数给出
$q_k\le C_{\rm sh}K_R$，其中 $C_{\rm sh}>0$ 固定。因此
当 $K_R<C_{\rm sh}^{-1}$ 时非零频率范围为空。在
非空范围上，初等的格计数给出
\[
 \#\{k:0<q_k\le C_{\rm sh}K_R\}
       \ll_{\rm fixed}K_R+\sqrt{K_R}+1\ll_{\rm fixed}K_R.
\]
因此在这些二进段上 $C,d,k$ 有 $O(C_0D_0K_R)$ 种选择。把
方程~\eqref{eq:gram-periodic-poisson} 乘以这个计数以及
公共前置因子，至多得到
\[
 (1+|\nu|)^{J_{\mathrm{Gr}}} R^{-B}\frac Q{Y'}\frac{RP_a}{C_0D_0}
       \min\left\{1,\left(\frac{RP_aC_0D_0}{Y'}\right)^A\right\}.
\]
对 $\Lambda=Y'/(RP_a)$，所需的二进正和为
\[
 \sum_{i,j\ge0}2^{-i-j}\min\{1,(2^{i+j}/\Lambda)^A\}
 =\sum_{n\ge0}(n+1)2^{-n}\min\{1,(2^n/\Lambda)^A\}
 \ll_A\frac{1+\log(2+\Lambda)}{\Lambda}.
\]
当 $\Lambda\ge1$ 时，在 $2^n\asymp\Lambda$ 处把两个几何尾部
拆开；当 $\Lambda<1$ 时，左边有界，而
显式给出的右边大于一个正常数。
因此非例外频率贡献
\[
 \ll (1+|\nu|)^{J_{\mathrm{Gr}}}\frac Q{Y'}\frac{P_a^2}{Y'}Z^\epsilon
\]
在求和 $R$ 之后，选取固定的 $B>2$。

剩下的是计数例外非零频率。它们满足
$(k)=a_1^6e$，其中 $e$ 无六次幂因子，且支撑在 $C$ 的
素因子与固定支撑上。可能的 $e$ 至多有
$6^{\omega(C)+O_S(1)}\ll_\epsilon q_C^\epsilon$ 个。
对其中每个，壳层界 $q_k\le C_{\rm sh}K_R$ 只允许
$q_{a_1}\le(C_{\rm sh}K_R/q_e)^{1/6}$。若该范围非空，
其上界至少为一，理想计数给出
$O_{\rm fixed}(K_R^{1/6})$ 种选择；若 $q_e>C_{\rm sh}K_R$ 则
它为空。单位因子只付出一个固定因子。因此这样
的元素 $k$ 有 $O_\epsilon(K_R^{1/6}q_C^\epsilon)$ 个，包括
$K_R$ 的有界次单位范围。

使用 $K_R=RP_a/C_0$、$|L_C|\le q_C$ 以及平凡的
两列点计数。在乘以 $Q/Y'^3$ 之前，这些二进段上的
贡献为
\[
 \ll (1+|\nu|)^{J_{\mathrm{Gr}}} R^{-B}Z^\epsilon
       Y'^2(RP_a)^{1/6}C_0^{-1/6}D_0^{-1}.
\]
在缩减任意小幂次之后，关于 $C_0,D_0$ 的求和收敛，
且对同一个固定的 $B>2$，关于 $R$ 的求和也收敛。这
给出项 $(Q/Y')P_a^{1/6}Z^\epsilon$。加上对角项
与非例外项即证明了该命题。
\end{proof}

\subsection{下界的完成}

现在把前面两个估计应用于精确低分离格式中同一个
重标度和式的两个因子。

\begin{proposition}[补偿的低估计]\label{prop:probe-low}
对于方程~\eqref{eq:compensated-probe-definition} 中的探针、
方程~\eqref{old-eq:5.1} 中的几何构型以及每个
$\epsilon>0$，
\begin{equation}
 |I_{\eta,\mathrm{modified}}(Z)|
       \ll Z^{l_x/2+b/12+\epsilon}=Z^{3/16+\epsilon}.
 \label{old-eq:5.10}
\end{equation}
该指数与目标及槽网格无关；常数与下阈值可以
依赖于固定的算术数据、槽系统以及光滑检验函数。
\end{proposition}

\begin{proof}
对于固定的重标度子集与元组，令 $Q=q_{b_*}X'Y'$。
其分离为方程~\eqref{eq:low-separated}，其中含
$A_{m,\sigma,v}(Y')$ 与 $B^J_{m,\sigma}(Z)$。令
$r_J=q_{p_J}/Z^d$；固定的环形支撑给出
$0<c_J\le r_J\le C_J<\infty$，且对元组一致成立。确切地说，
\[
 \begin{gathered}
 X'=Z^{l_x-d}/r_J,\qquad Y'=Z^{l_y-d}/r_J,\\
 Q=q_{b_*}Z^{M'}/r_J^2,\qquad
 P_a=Y'^2/Q=q_{b_*}^{-1}Z^{1/8}.
 \end{gathered}
\]
特别地，当 $Z\ge q_{b_*}^8$ 时 $P_a\ge1$。放大所允许的
固定数据下阈值，使得 $Q,Y'\ge1$ 对元组比值也一致
成立；这是可能的，因为 $M'\ge1/2$ 且
$l_y-d\ge5/16$。剩余的长度不等式为
\[
 l_x-d\ge3/16,\qquad l_y-d\ge5/16,
 \qquad l_y-d-11b/6\ge1/12.
\]
最后一个给出
\[
 \frac{Y'}{P_a^{11/6}}
 =q_{b_*}^{11/6}r_J^{-1}Z^{l_y-d-11b/6}
       \gg_{\rm fixed}Z^{1/12},
\]
从而 $P_a^2/Y'\ll P_a^{1/6}$。因此命题~\ref{prop:probe-gram}
给出
\begin{equation}
 \sum_{q_m\ll Q}|A_{m,\sigma,v}(Y')|^2
 \ll (1+|v|)^{J_{\mathrm{Gr}}}(Q/Y')P_a^{1/6}Z^\epsilon.
 \label{old-eq:5.6}
\end{equation}
在行和上应用 Cauchy--Schwarz、引理~\ref{lem:probe-row-norm}，
以及 $(1+|v|)^{J_{\mathrm{Gr}}/2}|\widehat W_0(iv)|$ 的可积性，现在
把该元组的未重标度分离界定在 $Z^\epsilon$ 乘以
\[
 \begin{aligned}
 Q^{-1/2}\bigl[(Q/Y')P_a^{1/6}\bigr]^{1/2}
       \bigl[Z^{M'}\bigr]^{1/2}
 &=\bigl[Z^{M'}/Y'\bigr]^{1/2}P_a^{1/12}\\
 &=r_J(X')^{1/2}P_a^{1/12}.
 \end{aligned}
\]
由于 $r_J$ 有界，这等于
$O((X')^{1/2}P_a^{1/12}Z^\epsilon)$。

对 $d=\sum_{i\in J}\ell_i$，重标度元组至多有 $Z^{d+\epsilon}$ 个，
系数为 $O(Z^{-3d/2})$，且
$(X')^{1/2}\asymp X^{1/2}Z^{-d/2}$。因此总的
附加指数为
\begin{equation}
 f(d)=d-3d/2-d/2=-d\le0.
 \label{old-eq:5.5}
\end{equation}
选取各分量估计中的小幂次，使其总和落在
指定的 $\epsilon$ 之内。对有限多个子集求和即证明
方程~\eqref{old-eq:5.10}。对于方程~\eqref{eq:part-II-mellin-target} 中
已定义的 $C=C_{\mathrm{II}}$，用于主信号的
正规化满足
\begin{equation}
 \begin{gathered}
 C_{\mathrm{II}}(s)=l_x/2+s-1+h/6=s-11/16,\\
 C_{\mathrm{II}}(7/8)=3/16=l_x/2+b/12.
 \end{gathered}
 \label{old-eq:5.11}
\end{equation}
因此，在第 II 部分所假设的矛盾之下，这个下界比
$Z^{C_{\mathrm{II}}(\beta_*)}$ 小一个精确幂次
$\Delta=\beta_*-7/8$，至多相差任意指定的 $\epsilon$。
\end{proof}

\section{局部补偿及其误差}
\label{sec:local-compensation}

节~\ref{sec:compensation} 中的物理修正已经
由其分离的低表示作了估计。我们现在确定其
完整的 Euler 修正，并界定两式运算所留下的局部误差。
全文之中，局部记号 $P_p,P_p^*,H_p,V,W,D,R$ 沿用
引理~\ref{lem:local-euler} 中的含义；在共享的解析估计中，其变量
$x$ 记作 $s$。固定一个非零且无六次幂因子的物理行 $u$，满足
$(u,S)=1$。所有槽素因子都属于节~\ref{sec:compensation}
的集合 $\mathcal P_i(Z)$；对每个这样的素因子令 $Q=q_p$。如
引理~\ref{lem:local-euler} 那样，记 $x_r=\Re x$、$w_r=\Re w$、$z_r=\Re z$。
$P_p^*$ 中的因子 $W_{\rm loc}=\rho Q^{-w}$ 与
$W=\chi_p(u)Q^{-w}$ 不同：当 $p\mid u$ 时，后者与 $D$ 都为零，而
前者保留其单位相位。

\subsection{完整的全纯修正}

对固定的槽素因子 $p$，把完备化指标限制在
$p\mid A$ 会在高次级数中把 $P_p$ 换成 $P_p^*$。把
$Z$ 变为 $ZQ$ 会使其 Mellin 权乘以 $Q^{x+z-1}$。因此
在 $P_p$ 有定义且非零的点处，方程~\eqref{eq:compensated-probe-definition}
的标记项具有乘子
$\overline{\eta(p)}Q^{x+z-1}P_p^*/P_p$。重标度项具有乘子
\[
 Q^{-3/2}Q^{-(1/2-z)}Q^{-(w-1)}=Q^{z-w-1}.
\]
因此，在同一轨迹上，两式运算的精确局部乘子为
\begin{equation}\label{old-eq:5.13}
 Q^{z-1}\mathcal B_p,\qquad
 \mathcal B_p=\overline{\eta(p)}Q^xP_p^*/P_p-Q^{-w}.
\end{equation}
把联合局部替换定义为
\begin{equation}\label{old-eq:5.13c}
 G_p=
 \frac{(\overline{\eta(p)}Q^x-Q^{-w})(1-V)(1-W)P_p^*
       -Q^{-w}(1-VW)}{1-D}.
\end{equation}
在 Euler 区域内使方程~\eqref{old-eq:5.13} 中原始商有定义的
点处，有 $G_p=H_p\mathcal B_p$。
所显示的公式在该商无定义处也定义了
$G_p$。右边既没有 $P_p$
分母，也没有 $1-W$ 分母。引理~\ref{lem:local-euler} 中的界 $|R|,|V|,|D|<1$
连同其关于 $P_p^*$ 的公式表明，
$G_p$ 在两个所述的 Euler 区域内全纯，包括在 $P_p$ 或 $H_p$
的零点处。

在 Euler 区域内使所选商有定义的点处，
把完整的槽
乘子记作
\[
 \mathcal B_i(u;x,w,z)=\sum_{p\in\mathcal P_i(Z)}
           W_i(q_p/P_i)q_p^{z-1}\mathcal B_p.
\]
用于移动围道的修正是在不使用这些商的情况下定义的：
\begin{equation}\label{eq:holomorphic-selected-tuple}
 \mathfrak H_{\eta,u,Z}(x,w,z)=
 \sum_{(p_i)\in\prod_i\mathcal P_i(Z)}
 \prod_{i=1}^K\bigl[W_i(q_{p_i}/P_i)q_{p_i}^{z-1}G_{p_i}\bigr]
 \prod_{\substack{p\notin S\\p\notin\{p_1,\ldots,p_K\}}}H_p.
\end{equation}
这是在提取出方程~\eqref{eq:probe-high} 中标量商之后的完整修正。为验证这一断言，
固定一个允许的元组并令 $\mathcal P=\{p_1,\ldots,p_K\}$。在
绝对起始线上，所选运算把 $P_p$ 换成
\begin{equation}\label{eq:selected-local-operation}
 \overline{\eta(p)}Q^{x+z-1}P_p^*-Q^{z-w-1}P_p
 =\frac{1-D}{(1-V)(1-W)}Q^{z-1}G_p.
\end{equation}
该等式来自把
$P_p-P_p^*=(1-V)^{-1}+W/(1-W)$ 代入
方程~\eqref{old-eq:5.13c}。它提取出与在未被选中的素因子处相同的标量局部因子
$(1-V)^{-1}(1-W)^{-1}(1-D)$。每个被选中的
素因子恰出现一次，因为槽
支撑互不相交。因此该元组的完备因子是
方程~\eqref{eq:probe-high} 中的标量商乘以方程
~\eqref{eq:holomorphic-selected-tuple} 中它对应的被加项。这一推理不涉及
除以 $P_p$ 或 $H_p$。在 Euler 区域内使所有
原始被选商都有定义的点处，同一个有限和也分解为
$\mathcal H_{\eta,u}\prod_i\mathcal B_i$。

每个被加项中未被选中的乘积在两个 Euler
区域内都正规收敛。更定量地，在方程
~\eqref{eq:local-region-one} 的任意固定子区域中，令
$\epsilon_H=\min(\epsilon_0,1/50)>0$。引理
~\ref{lem:local-euler} 中逐素因子的缺陷界给出，且对所选元组一致成立，
\begin{equation}\label{eq:unselected-local-product}
 \prod_{\substack{p\notin S\\p\notin\mathcal P}}|H_p|
 \le
 \prod_{\substack{p\notin S\\p\nmid u}}
       (1+Cq_p^{-1-\epsilon_H})
 \prod_{p\mid u}(1+Cq_p^{-\epsilon_H})
 \ll_\epsilon q_u^\epsilon.
\end{equation}
第一个正乘积由理想计数收敛，第二个
满足除数乘积界。略去被选中的因子只会从这些
正控制项中去掉因子。在第二个 Euler 区域中，同样的
论证使用相应的缺陷指数 $-363/200$ 与 $-33/40$。
因此方程~\eqref{eq:unselected-local-product} 在那里也成立，
并有相应的正控制项。这些上界并不
断言 $H_p$ 非零。

对每个固定的 $Z$，方程~\eqref{eq:holomorphic-selected-tuple} 是
若干个乘积的有限和，这些乘积由全纯的被选因子与正规
收敛的未被选乘积组成。因此它在两个 Euler 区域内每一点的
某个邻域上全纯。它还满足
定义~\ref{def:stage-high-data} 中的全高度要求。
事实上，在任一区域内的固定实盒上，$P_p^*$ 的公式以及
有界地远离零的分母给出 $|G_p|\ll Q^B$，其中
$B$ 固定，且对所有虚部与单位相位一致成立。一个槽中有
$O(P_i)$ 个理想，且在那里 $q_p\asymp P_i$。因此有限元组
和与方程~\eqref{eq:unselected-local-product} 给出
\begin{equation}\label{eq:compensated-all-height-majorant}
 |\mathfrak H_{\eta,u,Z}(x,w,z)|
 \ll_\epsilon q_u^\epsilon\prod_iP_i^{z_r+B}
 \ll_\epsilon q_u^\epsilon Z^{B'}.
\end{equation}
其中 $B'$ 在该盒上固定。对 $q_u\le Z^D$，这就是
方程~\eqref{eq:stage-all-height-majorant}，即便高度指数
为零。这里的所有常数都在以后任何分部积分阶数之前
固定。

现在把精确有限运算与
方程~\eqref{eq:probe-poisson-identity} 及 \eqref{eq:probe-high}
结合，给出
\begin{equation}\label{eq:compensated-poisson-identity}
\begin{split}
 I_{\eta,\mathrm{modified}}(Z)
 ={}&\frac1{(2\pi i)^3}\int_{(3)}\int_{(3)}\int_{(2)}
       \mathcal W(X,Y,Z;x,w,z)\\
 &\quad\cdot\sum_u^{(6)}q_u^{-z}\overline{\xi(u)}
       \frac{\zeta_F^S(6z)L^S(w,\chi_\bullet(u))}
       {L^S(x,\eta\overline{\chi_\bullet(u)})}
       \mathfrak H_{\eta,u,Z}(x,w,z)\,dz\,dw\,dx.
\end{split}
\end{equation}
该恒等式在所示直线上绝对收敛，因为对
固定的 $Z$，它由有限多个绝对收敛的重标度
基础探针恒等式得到。它没有额外的外部 Euler 因子。它
对物理表达式 $I_{\eta,\mathrm{modified}}$ 验证了定义~\ref{def:stage-high-data}
的精确高表示，其几何构型为方程~\eqref{old-eq:5.1}。特别地，
$l_x/2-1+h/6=17/96-1+13/96=-11/16$，正如
方程~\eqref{eq:part-II-mellin-target} 所要求。它正是
命题~\ref{prop:probe-low} 中所界定的那个探针的恒等式。每个延拓
都使用方程~\eqref{eq:holomorphic-selected-tuple} 中的完整修正，
而不是各个商的某个整体定义乘积。

\subsection{动态局部误差}

在恒等射线素因子处，来自 $T$ 的每个固定相位都等于一，尽管
$\eta(p)$ 可以是任意单位。先假设 $p\nmid u$ 并令
$v=\chi_p(u)$。则 $D=\eta(p)v^{-1}Q^{-x}$，且
$\overline{\eta(p)}Q^xD=v^{-1}$。利用
$v^{-1}W=Q^{-w}$ 与 $\mathcal E_p=P_p^*+D$ 可得到精确的
正规化相消
\begin{equation}\label{old-eq:5.13b}
\begin{split}
 G_p+v^{-1}H_p={}&
 \frac{(1-V)(1-W)}{1-D}
 \Bigl[(v^{-1}-Q^{-w})\{V/(1-V)-D\}\\
 &\hspace{17mm}+(\overline{\eta(p)}Q^x+v^{-1}-Q^{-w})
                    \mathcal E_p\Bigr].
\end{split}
\end{equation}
这是一个关于 Euler 区域内全纯局部表达式的恒等式。
在这些区域内，在原始商的轨迹上其左边为
$H_p(\mathcal B_p+v^{-1})$。
为在该轨迹上直接看出相消，在正规化之前把
$P_p^*=-D+\mathcal E_p$ 与
$P_p=(1-V)^{-1}+W/(1-W)+P_p^*$ 代入左边。
常数 $-v^{-1}$、$v^{-1}-Q^{-w}$ 与
$(v^{-1}-Q^{-w})W/(1-W)=Q^{-w}$ 之和为零。
所得有理恒等式只含分母 $1-R$、$1-V$
与 $1-D$，它们在 Euler 区域内非零；因此它给出
那里所显示的全纯恒等式。
对 $p\mid u$，主项由原来的零延拓为零，
方程~\eqref{old-eq:5.13c} 反而变为
\[
 G_p=
 \overline{\eta(p)}Q^x(1-V)\mathcal E_p-Q^{-w}H_p.
\]

我们现在给出对剩余误差槽的界。其解析
输入被明确陈述：反射的本原分子在
所保留的高度处必须很小。用当前的固定数据以及
$\Theta=\langle\eta,\widehat T\rangle$ 实例化节~\ref{sec:detector}，
并使用其零延拓表示 $\nu(n)\chi_n(u)^\varsigma$，其中
$\nu\in\Theta$ 且 $\varsigma\in\{-1,1\}$。分子
$\chi_\bullet(u)$ 及其共轭属于 $\mathcal X_u$，
固定乘子为 $1\in\Theta$；而分母
$\eta\overline{\chi_\bullet(u)}$ 属于 $\mathcal X_u$，固定
乘子为 $\eta\in\Theta$。因此那里的带缓冲估计，连同
引理~\ref{lem:deleted-euler-factors} 的逆删除因子界，
在保留高度处提供了下文所需的反射分子假设：对于
下文所用的点，$\Re(1-w)=a+6e$ 且
$|\Im(1-w)|\le T_1$，位于带缓冲的矩形内。恒等射线
限制使每个固定的 $T$ 相位在
槽素因子处都等于一，同时不改变物理行或任何非单位零点。

\begin{proposition}[动态局部误差与导子分配]
\label{prop:probe-errors}
设 $51/100\le a\le1$、$0<e\le10^{-3}$，并取
\[
 x_r=a+16e,\qquad w_r=1-a-6e,\qquad z_r=17/50.
\]
选择 $P_0$ 关于 $e$ 与固定数据充分大。
对于 $U,T_1\ge1$ 以及无六次幂因子、满足 $(u,S)=1$、
$q_u\asymp U$ 的行 $u$，设 $\psi_u^*$
为诱导 $\chi_\bullet(u)$ 的本原特征。设
$\mathcal W_u$ 为
$\{w:\Re w=w_r,\ |\Im w|\le T_1\}$ 的任意子集。假设 $\psi_u^*$
非主，且对 $w\in\mathcal W_u$ 一致成立，
\[
 |L(1-w,\overline{\psi_u^*})|\ll U^\epsilon
\]
其中幂次为任意指定的小量。
对于该行在所示动态区域上，用其前面的槽和定义 \(\mathcal B_i\)，
把每个 \(\mathcal B_p\) 解释为 \(G_p/H_p\)，
其中 \(G_p\) 来自方程~\eqref{old-eq:5.13c}。所述的
\(P_0\) 选取使那里对每个槽素因子都有 \(H_p\ne0\)，正如证明开头
所证。该命题只在这个动态区域上断言各别的延拓。
把一个槽的主部与误差部分别定义为
\[
 \mathcal Q_i(u;z)=
       -\sum_{p\in\mathcal P_i(Z)}W_i(q_p/P_i)q_p^{z-1}\overline{\chi_p(u)},
 \qquad
 \mathcal D_i(u)=\mathcal B_i(u;x,w,z)-\mathcal Q_i(u;z).
\]
则对槽的每个子集 $I$ 以及同样的点 $w\in\mathcal W_u$，
\begin{equation}\label{old-eq:5.13e}
 \left|L^S(w,\chi_\bullet(u))\prod_{i\in I}\mathcal D_i(u)\right|
 \ll U^{a-1/2+O(e)+\epsilon}(3+T_1)^C
       \prod_{i\in I}P_i^{z_r-1/2+O(e)+\epsilon}.
\end{equation}
主部保留物理行、所有零掩码以及
固定的恒等射线限制。特别地，其中心
正规化为
\[
 \mathcal Q_i(u;z)=-P_i^{z-1/2}
       \left(P_i^{-1/2}\sum_{p\in\mathcal P_i(Z)}
          \overline{\chi_p(u)}W_i(q_p/P_i)(q_p/P_i)^{z-1}\right).
\]
在同样的点 $w\in\mathcal W_u$ 处，且仅在这个动态区域上，
完整修正是有如下有限分解
\begin{equation}\label{eq:dynamic-compensated-decomposition}
 \mathfrak H_{\eta,u,Z}
 =\mathcal H_{\eta,u}\prod_{i=1}^K(\mathcal Q_i+\mathcal D_i)
 =\mathcal H_{\eta,u}
   \sum_{I\subseteq\{1,\ldots,K\}}
       \prod_{i\in I}\mathcal D_i\prod_{i\notin I}\mathcal Q_i.
\end{equation}
\end{proposition}

\begin{proof}
在这个区域中 $\vartheta=(-w_r)_+\le6e$。引理~\ref{lem:local-euler} 的
证明中的估计强化为：当 $p\nmid u$ 时
$H_p-1=O(Q^{-1-10e})$，当 $p\mid u$ 时
$H_p-1=O(Q^{-10e})$。增大固定的 $P_0$，使
所有这些素因子满足 $|H_p-1|<1/2$。方程~\eqref{old-eq:5.13c} 中的全纯
因子 $G_p$ 除以 $H_p$，
现在在整个区域上定义了 $\mathcal B_p$。当 $w_r<0$ 时有界地远离零的
是 $H_p$，而不是可能奇异的原始因子 $P_p$。

